% =====================================================================
%  Secao 3: O que é Backtracking?
% =====================================================================
\section{O que é Backtracking?}

Backtracking é uma técnica de busca utilizada para encontrar soluções em problemas que possuem múltiplas possibilidades.

Sua principal característica é construir a solução gradualmente. Em vez de tentar resolver o problema de uma única vez, o algoritmo adiciona uma escolha por etapa.

Sempre que uma escolha leva a uma situação inválida, o algoritmo retorna para o estado anterior e tenta outra alternativa.

O processo pode ser resumido em quatro etapas:

\begin{passos}[Quatro etapas]
\begin{enumerate}[leftmargin=1.4em]
  \item Fazer uma escolha.
  \item Verificar se a escolha continua válida.
  \item Avançar para a próxima etapa.
  \item Voltar atrás caso seja necessário.
\end{enumerate}
\end{passos}

Essa abordagem garante que todas as possibilidades relevantes sejam analisadas.

\figura{fig:ciclo}%
  {As quatro etapas do Backtracking formam um ciclo: escolher, validar, avançar e voltar atrás.}%
  {figuras/tikz/ciclo}
